% Notes of the Montreal seminar, 1970 (cote 129, typed part): archive cote 129, site page 29 (PDF page 30).
% typescript. Draft compared: none.
% AI-transcribed by Claude Opus 5.5 (claude-opus-5-5, adaptive thinking, effort high), Message Batches API, batch mtl129_full (msgbatch_01AGnDKHwy6ATCM7ehUa6263), 2026-11-13; not yet checked by eye.
\scanpage{129/29}
\mspage{\hand{II.1}}

\uline{GROUPES FINIS LOCALEMENT LIBRES ET THEORIE CLASSIQUE DE DIEUDONNE}

1. \uline{Rappels généraux sur les groupes finis localement libres}

1.1. Soit $S$ un schéma et soit $G$ un schéma en groupes fini et
localement libre sur $S$ . Il est équivalent de dire que $G = \mathrm{Spec}\, A$ ,
où $A$ est une $\underline{O}_S$-algèbre quasi-cohérente finie localement libre, la
structure de groupe étant décrite par des homomorphismes de $\underline{O}_S$-algèbres
\[ \Delta : A \to A \otimes_{\underline{O}_S} A \]
\[ I : A \to A \]
\[ \varepsilon : A \to \underline{O}_S \]
où $\Delta$ , $I$ , $\varepsilon$ correspondent respectivement à la loi de composition, à
l'application $x \to x^{-1}$ et à la section unité, ces morphismes satisfai-%
sant des relations évidentes exprimant l'associativité, etc... Ainsi, $A$
muni de sa loi d'algèbre et de $\Delta$ devient une \uline{bigèbre}
% NOTE: a handwritten curved line runs from a caret between "algèbre" and "et" on the line above down around the following displayed formula, apparently meant to bring "mu : A (x) A -> A" in after "sa loi d'algèbre"
\margin{$\times$} \margin{$\checkmark$}
\[ \hand{\mu}\deleted{\illegible{1}} : A \otimes A \to A . \]

\margin{\hand{$\mathcal{J}$}} Le noyau \deleted{$J$}\inserted{$\mathcal{J}$} de $\varepsilon$ s'appelle \uline{l'idéal d'augmentation} du groupe $G$ .

\margin{\unsure{C.E.}} Dans tout ce qui suit, on supposera, sauf mention contraire,
que tous les groupes considérés sont \uline{commutatifs}.

1.2. Si $s \in S$ , on appelle \uline{rang de} $G$ \uline{en} $s$ le rang sur $k(s)$

\mspage{\hand{29}}
